\newpage
%%%%%%%%%%%%%%%%%%%%%%%%%%%%%%%%%%%%%%%%%%%%%%%%%%%%%%%%%%%%%%%%
%%%%%%%%%%%%%%%%%%%%%%%%%%%%%%%%%%%%%%%%%%%%%%%%%%%%%%%%%%%%%%%%
%%%%%%%%%%%%%%%%%%%%%%%%%% Enunciado %%%%%%%%%%%%%%%%%%%%%%%%%%%

\begin{myblock}
\phantomsection\addcontentsline{toc}{section}{Parte B | BETO}
\section*{Parte B | BETO}

    Usando la misma base de la tarea 1, entrena un modelo por cada arquitectura: 
    CNN, RNN/LSTM/GRU, y un Transformer de Hugging Face con fine-tuning. Reporta
    accuracy y F1, incluye matriz de confusión y análisis de errores característicos.  

\end{myblock}

Para esta prueba reutilizamos lo usado en RESTMex2025. 

Es claro como un trasnformer mejora sustancialmente los resultados, superando por bastante a las arquitecturas de redes neuronales
recurrentes. La jerarquía ahora queda establecida de manera muy clara: Transformers >> BiLSTM > BiGRU > Bi-RNN > TextCNN.


\subsection{Visualización de Modelos}

% --- Figura: Matrices de Confusión para Transformers ---

\begin{figure}[h!]
    \centering
    \includegraphics[width=0.8\linewidth]{/home/cesar/Natural-Language-Processing-and-Computer-Vision/Generacion-y-Clasificacion-de-Texto/Reporte/Images/matriz_confusion_transformer_c.png}
    \caption{Matriz de Confusión | Polaridad Transformer.}
    \label{fig:Matriz_Transformer}
\end{figure}


\clearpage


